\documentclass[12pt,oneside]{book}

% ============================================================
% METADATOS
% ============================================================
\newcommand{\universidad}{Universidad de Buenos Aires}
\newcommand{\materia}{Derecho Civil --- Parte General}
\newcommand{\titulo}{La Capacidad Jurídica}
\newcommand{\autor}{Isaac Edgar Camacho Ocampo}
\newcommand{\anio}{2026}

% ============================================================
% PAQUETES
% ============================================================
\usepackage[utf8]{inputenc}
\usepackage[T1]{fontenc}
\usepackage[spanish,es-nodecimaldot]{babel}

\usepackage[
    top=2.5cm,
    bottom=2.5cm,
    left=3cm,
    right=2.5cm,
    headheight=15pt
]{geometry}

\usepackage{amsmath}
\usepackage{amssymb}
\usepackage{mathtools}

\usepackage{graphicx}
\usepackage{float}
\usepackage{caption}
\usepackage{subcaption}

\usepackage{booktabs}
\usepackage{array}
\usepackage{longtable}
\usepackage{tabularx}

\usepackage{enumitem}

\usepackage{xcolor}
\usepackage[most]{tcolorbox}

\usepackage{fancyhdr}
\usepackage{ragged2e}

\graphicspath{
    {./img/}
    {./graficos/}
    {./figuras/}
}

% ============================================================
% CONFIGURACIÓN
% ============================================================
\setcounter{tocdepth}{2}

\setlength{\parindent}{0pt}
\setlength{\parskip}{0.5em}

\setlist[itemize]{
    topsep=0.3em,
    itemsep=0.2em
}
\setlist[enumerate]{
    topsep=0.3em,
    itemsep=0.2em
}

\clubpenalty=10000
\widowpenalty=10000

\pagestyle{fancy}
\fancyhf{}
\fancyhead[L]{\nouppercase{\leftmark}}
\fancyhead[R]{\materia}
\fancyfoot[C]{\thepage}
\renewcommand{\headrulewidth}{0.4pt}
\renewcommand{\footrulewidth}{0pt}

\fancypagestyle{plain}{
    \fancyhf{}
    \fancyfoot[C]{\thepage}
    \renewcommand{\headrulewidth}{0pt}
}

% ============================================================
% COMANDOS MATEMÁTICOS
% ============================================================
\newcommand{\abs}[1]{\left|#1\right|}
\newcommand{\norm}[1]{\left\lVert#1\right\rVert}
\newcommand{\vect}[1]{\mathbf{#1}}
\newcommand{\deriv}[2]{\frac{d#1}{d#2}}
\newcommand{\pderiv}[2]{\frac{\partial#1}{\partial#2}}

% ============================================================
% ENTORNOS / CAJAS ACADÉMICAS
% ============================================================
\newtcolorbox{definition}[1]{
    enhanced,
    breakable,
    title={Definición: #1},
    fonttitle=\bfseries,
    colback=blue!3,
    colframe=blue!50!black,
    boxrule=0.8pt,
    arc=2mm
}

\newtcolorbox{important}{
    enhanced,
    breakable,
    title={Importante},
    fonttitle=\bfseries,
    colback=yellow!8,
    colframe=orange!70!black,
    boxrule=0.8pt,
    arc=2mm
}

\newtcolorbox{example}[1]{
    enhanced,
    breakable,
    title={Ejemplo: #1},
    fonttitle=\bfseries,
    colback=green!4,
    colframe=green!40!black,
    boxrule=0.8pt,
    arc=2mm
}

\newtcolorbox{formula}[1]{
    enhanced,
    breakable,
    title={Fórmula / Regla: #1},
    fonttitle=\bfseries,
    colback=purple!4,
    colframe=purple!50!black,
    boxrule=0.8pt,
    arc=2mm
}

\newtcolorbox{exercise}[1]{
    enhanced,
    breakable,
    title={Ejercicio: #1},
    fonttitle=\bfseries,
    colback=gray!5,
    colframe=gray!60!black,
    boxrule=0.8pt,
    arc=2mm
}

\newtcolorbox{note}{
    enhanced,
    breakable,
    title={Nota},
    fonttitle=\bfseries,
    colback=white,
    colframe=gray!50,
    boxrule=0.6pt,
    arc=2mm
}

% ============================================================
% CONFIGURACIÓN FINAL: HIPERVÍNCULOS Y METADATOS PDF
% ============================================================
\usepackage[
    colorlinks=true,
    linkcolor=blue!50!black,
    citecolor=green!40!black,
    urlcolor=blue!60!black,
    pdftitle={\titulo},
    pdfauthor={\autor},
    pdfsubject={Apuntes universitarios},
    bookmarks=true
]{hyperref}

% ============================================================
% DOCUMENTO
% ============================================================
\begin{document}

% ---------------- PORTADA ----------------
\begin{titlepage}
\thispagestyle{empty}

\begin{center}

\vspace*{1cm}

{\large \textsc{\universidad}}

\vspace{3cm}

{\Huge \textbf{\materia}}

\vspace{1cm}

{\LARGE \titulo}

\vspace{0.8cm}

\textit{Comprender los conceptos es el primer paso para convertir el estudio en conocimiento.}

\vspace{1.2cm}

\rule{0.70\textwidth}{0.5pt}

\vspace{0.5cm}

{\large Apuntes de estudio}

\vspace{0.5cm}

\rule{0.70\textwidth}{0.5pt}

\vfill

{\large \autor}

\vspace{0.3cm}

{\large \anio}

\end{center}

\vfill

\begin{flushleft}
\textit{Este es un modesto aporte a los estudiantes de ``Derecho Civil'' no pretende sustituir el material oficial de la materia.}
\end{flushleft}

\end{titlepage}

% ---------------- ÍNDICE ----------------
\tableofcontents
\clearpage

% ============================================================
% CONTENIDO
% ============================================================
\chapter{La Capacidad Jurídica}

Para comprender esta unidad resulta útil seguir un caso guía que atraviesa todos los temas: una niña de corta edad hereda un inmueble de su padre fallecido. Ella es titular del derecho sobre el bien (capacidad de derecho), pero no puede ejercerlo por sí misma (carece de capacidad de ejercicio), por lo que sus padres o un tutor la representan. Si, siendo ya adulta, un juez determinara que necesita apoyo para administrar bienes, se le designaría un asistente y no un representante. En todo ese recorrido, el Ministerio Público vela por sus intereses. Este hilo conductor permite relacionar los distintos bloques temáticos que se desarrollan a continuación.

\begin{note}
La capacidad jurídica constituye el punto de partida del Derecho Civil: determina quién puede ser titular de derechos y quién puede ejercerlos por sí mismo. Esta cuestión resulta relevante en la celebración de contratos, en los procesos sucesorios, en la representación judicial y en la protección de personas menores de edad o con discapacidad.
\end{note}

\section{Concepto de capacidad jurídica}

\subsection{Definición de capacidad}

La capacidad es una cierta aptitud —una aptitud para entablar relaciones jurídicas— que presenta dos componentes: un componente de \textbf{derecho} (aptitud para ser titular de relaciones jurídicas, de derechos y obligaciones) y un componente de \textbf{hecho o de ejercicio} (aptitud para ejercer esos derechos). Ambas aptitudes se regulan bajo un concepto común en el Código: el de \textbf{capacidad jurídica}.

\subsection{Capacidad jurídica como concepto unificador}

La capacidad jurídica es un constructo común que engloba la capacidad de derecho y la capacidad de ejercicio —esta última antes también llamada capacidad de hecho—. El tratamiento del tema parte de una regla general: la regla de la capacidad.

\begin{definition}{Regla general de la capacidad (art. 22 CCyCN)}
Toda persona humana goza de la aptitud para ser titular de derechos y deberes jurídicos. La ley puede privar o limitar esa capacidad respecto de hechos, simples actos o actos jurídicos determinados.
\end{definition}

La razón de fondo es existencial: si los seres humanos se desarrollan entablando relaciones jurídicas, negar esa aptitud equivaldría a impedirles satisfacer sus necesidades básicas y desplegar su personalidad para el logro de sus fines. Ese despliegue requiere la aptitud de ser titular de derechos y deberes, comenzando por los derechos humanos, que no son concedidos por el ordenamiento jurídico sino reconocidos, pues le corresponden al ser humano por el solo hecho de serlo: el derecho a la libertad, la libertad de contrato, la libertad de asociación y otras libertades básicas.

\subsection{Capacidad de derecho}

Algunos autores señalan que la capacidad de derecho se identifica con la propia condición de persona: se es persona y, por eso, se tiene capacidad de derecho. Un ser humano con incapacidad absoluta de derecho equivaldría a un muerto civil, alguien sin posibilidad alguna de desarrollarse.

Por ello, la capacidad de derecho se reconoce a toda persona. Pueden existir algunas incapacidades, limitaciones o privaciones —que la tradición clásica denomina incapacidades de derecho relativas— para ciertas circunstancias en las que se priva o limita esa capacidad.

\begin{important}
Las incapacidades de derecho son siempre \textbf{relativas} (nunca absolutas), se aplican a \textbf{actos puntuales} y están dispuestas para proteger el \textbf{interés general}. No pueden ser suplidas por medio de un representante y su violación genera \textbf{nulidad absoluta}.
\end{important}

\subsection{Capacidad de ejercicio (o de hecho)}

La capacidad de ejercicio es la aptitud para desplegar, para gozar de los derechos, por sí mismo. Allí radica la clave del concepto: la actuación por sí mismo.

A diferencia de las limitaciones a la capacidad de derecho, las limitaciones a la capacidad de ejercicio están puestas para \textbf{proteger a la propia persona}.

\begin{example}{Incapacidad de ejercicio de un menor de edad}
A una persona menor de edad se la considera incapaz de hecho o de ejercicio no porque carezca de derechos, sino porque, en razón de su inmadurez, se le proveen representantes y una serie de salvaguardias que permiten tutelarla por su vulnerabilidad.
\end{example}

\section{Incapacidades de derecho: el art. 1002 CCyCN}

El ejemplo más ilustrativo de incapacidades de derecho lo brinda el artículo 1002, referido a las inhabilidades para contratar.

\begin{formula}{Inhabilidades para contratar (art. 1002 CCyCN)}
No pueden contratar en interés propio:
\begin{itemize}
    \item Los funcionarios públicos respecto de bienes cuya administración o enajenación está a su cargo.
    \item Los jueces, funcionarios y auxiliares de la justicia, árbitros, mediadores y sus auxiliares, respecto de bienes relacionados con procesos en los que intervienen o han intervenido.
    \item Los abogados y procuradores respecto de bienes litigiosos en procesos en los que intervienen o han intervenido.
    \item Los cónyuges bajo el régimen de comunidad entre sí.
    \item Los albaceas que no son herederos respecto de bienes de las testamentarías que están a su cargo.
\end{itemize}
\end{formula}

La razón de esta regla es de interés público: no se admite que el funcionario que administra un bien sea quien aparezca comprándolo, que un árbitro que interviene en un conflicto entre particulares compre los bienes en disputa, o que un juez adquiera la casa que está en remate en su propio juzgado, pues ello configuraría un conflicto de intereses.

\begin{table}[H]
\centering
\caption{Incapacidad de derecho vs. incapacidad de ejercicio}
\begin{tabularx}{\textwidth}{>{\raggedright\arraybackslash}p{0.28\textwidth} X X}
\toprule
\textbf{Característica} & \textbf{Incapacidad de derecho} & \textbf{Incapacidad de ejercicio} \\
\midrule
¿A quién protege? & Al interés general / público & A la propia persona \\
¿Puede ser absoluta? & Nunca (siempre relativa) & Sí (casos extremos: art. 32) \\
¿Se suple con representante? & No & Sí \\
¿Alcance? & Actos puntuales determinados & Puede ser amplio \\
¿Regulación? & A lo largo de todo el Código & Art. 24 y ss. CCyCN \\
Sanción por violación & Nulidad absoluta & Acto anulable (según el caso) \\
\bottomrule
\end{tabularx}
\end{table}

\section{Incapaces de ejercicio}

El Código presenta distintos supuestos de quiénes son incapaces de ejercicio.

\begin{definition}{Incapaces de ejercicio (art. 24 CCyCN)}
Son incapaces de ejercicio:
\begin{enumerate}[label=(\alph*)]
    \item La persona por nacer.
    \item La persona que no cuenta con la edad y grado de madurez suficiente, con el alcance dispuesto en la Sección 2ª de este Capítulo.
    \item La persona declarada incapaz por sentencia judicial, en la extensión dispuesta en esa decisión.
\end{enumerate}
\end{definition}

\subsection{Personas por nacer}

La persona por nacer, conforme al artículo 19, es aquella que está concebida hasta el momento del nacimiento. Son personas, cuentan con la representación de sus padres (art. 101) y pueden recibir bienes de tipo patrimonial, pero la adquisición de derechos patrimoniales queda sujeta al nacimiento con vida (art. 21).

\subsection{Personas menores de edad}

Una vez nacida, la persona menor de edad —definida en el artículo 25 como quien no cuenta con la edad y grado de madurez suficiente— es, por regla, incapaz de ejercicio, pero va adquiriendo una capacidad progresiva a medida que crece en edad y madurez.

\begin{important}
Desde que nace hasta que cumple 18 años —edad en que se alcanza la mayoría de edad—, la persona va evolucionando. La representación que corresponde en primer lugar a sus padres (responsabilidad parental) va decreciendo a medida que las aptitudes de la persona menor crecen, hasta que se hace plenamente capaz.
\end{important}

El Estado reconoce a los padres una autoridad para educar porque los considera en las mejores condiciones para criar a sus hijos. Puede suceder algo que los prive de la responsabilidad parental en casos muy graves, pero la regla es que son los padres quienes van representando al hijo que han traído a la vida.

\subsection{Personas declaradas incapaces por sentencia judicial}

Son los supuestos del artículo 32, último párrafo. Se trata de casos excepcionales en los que la persona está absolutamente imposibilitada de interaccionar con su entorno, no puede expresar su voluntad de ninguna manera ni por ningún medio, y el sistema de apoyo resulta ineficaz.

\begin{definition}{Declaración de incapacidad (art. 32, último párrafo, CCyCN)}
Por excepción, cuando la persona se encuentre absolutamente imposibilitada de interaccionar con su entorno y expresar su voluntad por cualquier modo, medio o formato adecuado, y el sistema de apoyos resulte ineficaz, el juez puede declarar la incapacidad y designar un curador.
\end{definition}

En estos casos, el juez debe determinar en qué actos y funciones puede actuar la persona, y se le nombra un curador. Son casos muy extremos, generalmente de personas con una discapacidad severa que no pueden interaccionar de ninguna manera.

\begin{example}{Caso de Juana Fernández}
Juana Fernández tiene 3 años. Su padre fallece y ella resulta única heredera. La herencia comprende un inmueble y dinero en una cuenta bancaria. Hay que tomar una serie de decisiones sobre esos bienes, pero no las va a tomar una niña de tres años, por lo que se nombra un tutor o representante que se ocupará de ejercer lo que ella no puede hacer por sí. Los bienes están a nombre de Juana Fernández (capacidad de derecho: los bienes son de ella), pero ella no puede ejercer ese derecho por sí misma a los tres años; el tutor interviene precisamente en su capacidad de ejercicio.
\end{example}

\section{Capacidad restringida y sistemas de apoyo}

\subsection{Personas con capacidad restringida}

Existe otro mecanismo cuando una persona tiene capacidad restringida, es decir, cuando es capaz de ejercicio pero una sentencia judicial establece que ciertos actos no puede realizarlos por sí sola, sino que necesita un apoyo.

Esto constituye, en la actualidad, prácticamente la regla en el sistema del Código Civil y Comercial, especialmente a partir de la Convención sobre los Derechos de las Personas con Discapacidad, en virtud de la cual se considera que las personas con discapacidad deben ser integradas a la sociedad y pueden expresar sus preferencias aun cuando necesiten apoyos.

\begin{example}{Restricción de la capacidad para actos patrimoniales}
Una persona con discapacidad mental quizás puede expresarse y manejarse con pequeñas sumas de dinero, pero no comprende el valor de un inmueble. El juez, en un proceso judicial, restringe su capacidad: para las decisiones de la vida cotidiana actúa libremente, pero para la venta de inmuebles necesita el apoyo de un familiar.
\end{example}

\subsection{Sistema de apoyos: diferencia con la representación}

En el supuesto de capacidad restringida no se está ante un caso de representación. El apoyo no ocupa el lugar de la persona con discapacidad ni ejerce los derechos en su nombre. Se trata de un caso de \textbf{asistencia}, porque concurren ambas voluntades: la de la persona con capacidad restringida y la del apoyo, con las funciones que determina el juez.

\begin{note}
Si el apoyo no da su consentimiento —por ejemplo, porque advierte que se está haciendo un mal negocio, vendiendo muy barato en un momento de caída de precios— puede negarse. El apoyo existe para brindar seguridad, para aconsejar y para evitar que se aproveche a la persona que quizás no comprende bien el valor del dinero.
\end{note}

\begin{table}[H]
\centering
\caption{Representación vs. asistencia (sistema de apoyos)}
\begin{tabularx}{\textwidth}{>{\raggedright\arraybackslash}p{0.24\textwidth} X X}
\toprule
\textbf{Aspecto} & \textbf{Representación} & \textbf{Asistencia (apoyos)} \\
\midrule
¿Para quién? & Incapaces de ejercicio & Personas con capacidad restringida \\
¿Quién actúa? & El representante en nombre del incapaz & Ambos: persona + apoyo \\
¿La voluntad del titular? & No es necesaria & Es necesaria y principal \\
¿Quién designa? & La ley o el juez & El juez (proceso específico) \\
Finalidad & Sustituir al incapaz & Acompañar y proteger \\
\bottomrule
\end{tabularx}
\end{table}

\subsection{Personas inhabilitadas (pródigos)}

\begin{definition}{Inhabilitación por prodigalidad (art. 48 CCyCN)}
Pueden ser inhabilitados quienes, por la prodigalidad en la gestión de sus bienes, expongan a su cónyuge, conviviente o a sus hijos menores de edad o con discapacidad a la pérdida del patrimonio. A estos efectos, se considera que hay prodigalidad cuando, sin mediar razones atendibles, se dilapidan los bienes.
\end{definition}

El pródigo es quien dilapida su patrimonio, poniendo en riesgo también a su familia. En estos casos, judicialmente se pueden tomar medidas: si se inicia una causa por prodigalidad, se restringe la capacidad de la persona y se le nombran apoyos para la administración de sus bienes.

\subsection{Representantes legales (art. 101 CCyCN)}

\begin{definition}{Representantes legales (art. 101 CCyCN)}
Son representantes:
\begin{enumerate}[label=(\alph*)]
    \item De las personas por nacer, sus padres.
    \item De las personas menores de edad no emancipadas, sus padres. Si faltan los padres, o son incapaces, o han sido privados de la responsabilidad parental, un tutor.
    \item De las personas con capacidad restringida y de las personas declaradas incapaces, los apoyos (cuando tengan representación) y el curador, respectivamente.
\end{enumerate}
\end{definition}

\section{Ministerio Público (art. 103 CCyCN)}

El Ministerio Público es una institución con fundamento constitucional, ubicada como un organismo extrapoder, que en materia civil actúa velando por los intereses de las personas menores de edad e incapaces o con capacidad restringida. Está dividido en Ministerio Público de la Defensa y Ministerio Público Fiscal; en esta materia el análisis se concentra en el Ministerio Público de la Defensa.

\begin{note}
Es una forma complementaria de protección: a los incapaces de ejercicio se les provee un representante o un asistente según el caso, pero además, en todos los casos, el Ministerio Público vela por sus intereses.
\end{note}

\begin{definition}{Actuación del Ministerio Público (art. 103 CCyCN)}
La actuación del Ministerio Público respecto de personas menores de edad, incapaces y con capacidad restringida, y de aquellas cuyo ejercicio de capacidad requiera de un sistema de apoyos, puede ser, en el ámbito judicial, complementaria o principal.
\end{definition}

\subsection{Actuación judicial complementaria}

La actuación complementaria corresponde al Ministerio Público en todos los procesos en los que se encuentren involucradas personas menores de edad, incapaces o con capacidad restringida. Consiste en darle vista de las actuaciones.

\begin{example}{Vista al Ministerio Público en un proceso sucesorio}
Fallece uno de los progenitores. La persona menor de edad también hereda junto al otro progenitor, y se tramita un juicio de sucesión con varios herederos en una situación compleja. De cada paso judicial que toque los intereses del menor se dará vista al Ministerio Público —conocido como asesor de menores o defensor de menores e incapaces—, que intervendrá cada vez que esté en juego un interés de una persona menor o incapaz.
\end{example}

Esta intervención se produce sin sustituir al representante: si sobrevive la madre, ella actúa como representante principal, mientras que el Ministerio Público también dictamina velando por los intereses del menor, pero de manera complementaria, sin desplazar a la representación principal.

\subsection{Actuación judicial principal}

La actuación judicial principal se produce cuando los intereses de los representados están comprometidos y los representantes no accionan, cuando se trata de exigir a los propios representantes que cumplan con sus deberes, o cuando se carece de representante legal.

\begin{example}{Actuación principal del Ministerio Público}
\textbf{Caso 1:} una persona menor de edad tiene un solo progenitor que no le pasa alimentos; el Ministerio Público toma la representación y promueve un juicio de alimentos contra ese progenitor. \textbf{Caso 2:} una niña queda huérfana; se da intervención al Ministerio Público, que actúa ante el juez como si fuera el representante, en carácter de actuación judicial principal, sustituyendo a los representantes que menciona el artículo 101.
\end{example}

\subsection{Actuación extrajudicial}

La actuación extrajudicial está habilitada en forma extraordinaria en el artículo 103. En este supuesto no existe proceso judicial: el Ministerio Público interviene directamente para preservar los derechos económicos, sociales y culturales de menores o incapaces ante la ausencia, carencia o inacción de sus representantes legales.

\begin{example}{Intervención extrajudicial}
Un niño o niña está sometido a un ambiente de drogas, corrompido y peligroso, que pone en riesgo su situación, y los padres no hacen nada o son promotores de esa situación. El Ministerio Público toma cartas en el asunto e interviene de manera extrajudicial para preservar los derechos del menor.
\end{example}

% ============================================================
% CORNELL
% ============================================================
\section{Cuadro Cornell --- Repaso de la unidad}

\begin{longtable}{
    >{\raggedright\arraybackslash}p{0.28\textwidth}
    >{\raggedright\arraybackslash}p{0.62\textwidth}
}
\toprule
\textbf{Preguntas / palabras clave} & \textbf{Notas / respuestas} \\
\midrule
\endhead

\textbf{¿Qué es la capacidad jurídica?} &
Es la aptitud para entablar relaciones jurídicas. Engloba dos componentes: la capacidad de derecho (ser titular de derechos y deberes) y la capacidad de ejercicio (actuar por sí mismo). El Código Civil y Comercial (CCyCN) las regula conjuntamente bajo este concepto (art. 22). \\

\textbf{¿Cuál es la regla general del art. 22?} &
Toda persona humana goza de la aptitud para ser titular de derechos y deberes jurídicos. La ley solo puede privar o limitar esa capacidad respecto de hechos, simples actos o actos jurídicos determinados. La incapacidad es siempre la excepción. \\

\textbf{¿En qué se diferencia la capacidad de derecho de la de ejercicio?} &
La capacidad de derecho es la aptitud para ser titular de derechos (nunca puede ser absolutamente negada; su incapacidad es siempre relativa y para actos puntuales). La capacidad de ejercicio es la aptitud para actuar por sí mismo; puede restringirse para proteger a la propia persona y suplirse con representación o asistencia. \\

\textbf{¿Qué son las incapacidades de derecho relativas?} &
Son privaciones o limitaciones de la aptitud para ser titular de ciertos derechos en situaciones específicas, impuestas para proteger el interés general. Ejemplo paradigmático: el art. 1002 prohíbe a los funcionarios públicos contratar sobre bienes que administran. No pueden ser suplidas por representante y su violación genera nulidad absoluta. \\

\textbf{¿Quiénes son incapaces de ejercicio?} &
Según el art. 24 CCyCN: (a) la persona por nacer, (b) la persona menor de edad (hasta los 18 años, con capacidad progresiva) y (c) la persona declarada incapaz por sentencia judicial en casos excepcionales, cuando está absolutamente imposibilitada de interaccionar con su entorno (art. 32, último párrafo). \\

\textbf{¿Qué es la capacidad progresiva?} &
Principio por el cual la persona menor de edad va adquiriendo progresivamente mayor autonomía a medida que crece en edad y madurez. La representación parental decrece en la misma medida en que la capacidad del menor crece, hasta la mayoría de edad, a los 18 años. \\

\textbf{¿Qué es la capacidad restringida y cómo difiere de la incapacidad?} &
La capacidad restringida (art. 32) implica que la persona es capaz de ejercicio, pero una sentencia judicial determina que ciertos actos requieren apoyo. Es la regla hoy en día, influida por la Convención sobre Derechos de Personas con Discapacidad. La incapacidad declarada, en cambio, es la excepción extrema para quienes no pueden interaccionar de ningún modo. \\

\textbf{¿Qué diferencia hay entre representación y asistencia?} &
En la representación (para incapaces), el representante actúa en nombre del incapaz ejerciendo sus derechos. En la asistencia o apoyo (para personas con capacidad restringida), concurren ambas voluntades: la de la persona y la del apoyo. El apoyo no sustituye, sino que acompaña. \\

\textbf{¿Quiénes son los inhabilitados y qué los diferencia de los incapaces?} &
Los inhabilitados son personas que, por prodigalidad (dilapidar su patrimonio), ponen en riesgo a su familia (art. 48). Se les restringe la capacidad y se les nombran apoyos para la administración de bienes. A diferencia de los incapaces, no pierden el ejercicio de todos sus actos, sino solo los patrimoniales relevantes. \\

\textbf{¿Cuál es el rol del Ministerio Público en materia de capacidad?} &
Es una forma complementaria de protección que vela por los intereses de menores, incapaces y personas con capacidad restringida. Actúa judicialmente de forma complementaria (recibe vista en todos los procesos que afecten al menor) o principal (cuando los representantes no actúan o hay conflicto). También puede actuar extrajudicialmente en casos urgentes de vulneración de derechos económicos, sociales y culturales (art. 103). \\

\midrule

\multicolumn{2}{p{0.94\textwidth}}{
\textbf{Resumen:}
La capacidad jurídica es la aptitud que tienen las personas humanas para entablar relaciones jurídicas y desenvolverse en el ordenamiento civil. El Código Civil y Comercial de la Nación parte de una regla de capacidad (art. 22): toda persona es titular de derechos y deberes, y solo la ley puede establecer limitaciones para actos determinados. Esta regla se desdobla en dos dimensiones: la capacidad de derecho —ser titular de derechos—, que admite incapacidades relativas para proteger el interés general y que no puede ser suplida por representación (art. 1002 como ejemplo paradigmático); y la capacidad de ejercicio —actuar por sí mismo—, que puede limitarse para proteger a la propia persona. Son incapaces de ejercicio las personas por nacer, los menores de edad (con capacidad progresiva) y quienes son declarados incapaces por sentencia en casos extremos (art. 32). Para quienes tienen capacidad restringida —la regla en la actualidad, influida por la Convención sobre Derechos de Personas con Discapacidad—, el ordenamiento prevé sistemas de apoyo que asisten (no sustituyen) la voluntad de la persona. Frente a los incapaces se erige la representación; frente a los de capacidad restringida, la asistencia. En todos los casos, el Ministerio Público actúa como protección adicional: complementariamente en los procesos judiciales donde haya menores o incapaces, y de forma principal cuando los representantes legales no cumplen sus funciones o están ausentes (art. 103 CCyCN).
} \\

\bottomrule
\end{longtable}

\end{document}
